% Start the method example on a fresh page with its opening figure.
\newpage
\section{Problem Formulation and Method}
\label{sec:method}
\begin{figure}[htbp]
  \centering
  % Pure vector TikZ figure, reused at full text width in both layouts.
\begin{tikzpicture}[x=0.01\linewidth,y=1cm,
  stage/.style={draw=STARLine,fill=STARLight,rounded corners=2pt,
    minimum height=1.02cm,text width=0.185\linewidth,align=center,
    inner sep=5pt,font=\sffamily\small},
  flow/.style={-{Latex[length=2mm]},line width=0.65pt,draw=STARForest}]
  \node[stage] (a) at (11,0) {\textbf{Specify}\\[-1pt]
    {\scriptsize Data and assumptions}};
  \node[stage] (b) at (37,0) {\textbf{Construct}\\[-1pt]
    {\scriptsize Model and program}};
  \node[stage] (c) at (63,0) {\textbf{Execute}\\[-1pt]
    {\scriptsize Compiler and runtime}};
  \node[stage] (d) at (89,0) {\textbf{Validate}\\[-1pt]
    {\scriptsize Evidence and artifacts}};
  \draw[flow] (a) -- (b);
  \draw[flow] (b) -- (c);
  \draw[flow] (c) -- (d);
  \draw[STARForest,line width=0.65pt]
    (1,-0.80) -- (1,-0.98) -- (99,-0.98) -- (99,-0.80);
  \node[fill=white,inner xsep=7pt,font=\sffamily\scriptsize,
    text=STARForest] at (50,-0.98) {EXPLICIT CONTRACTS AT EVERY BOUNDARY};
\end{tikzpicture}

  \caption{An illustrative research pipeline. Each boundary carries an
  explicit contract: what enters, what changes, and what is checked.
  This diagram is a template component, not an evaluated system.}
  \label{fig:pipeline}
\end{figure}

\subsection{Define the contract}
Let $\dataset=\{(x_i,y_i)\}_{i=1}^{n}$ be a training set and let
$f_\theta$ denote a model with parameters $\theta\in\R^d$.
Write the objective in a form that separates task quality from the
system property under study:
\begin{equation}
  \theta^{\star}=\arg\min_{\theta}\;
  \frac{1}{n}\sum_{i=1}^{n}\loss(f_\theta(x_i),y_i)
  +\lambda\,\Omega(\theta;\mathcal{S}).
  \label{eq:objective}
\end{equation}
Here $\mathcal{S}$ specifies the deployment setting, $\Omega$ is a chosen
penalty, and $\lambda\geq 0$ controls its weight. Define the units and
estimation procedure for each term. A penalty is not, by itself, a guarantee
that a deployed system satisfies the intended property.

\begin{assumption}[Fixed evaluation interface]
\label{ass:interface}
The candidate and reference implementations use the same input schema,
preprocessing, and declared numerical precision during evaluation.
\end{assumption}

\subsection{Make the mechanism inspectable}
\autoref{fig:pipeline} separates specification, construction, execution,
and validation. Define each stage's inputs, transformation, and outputs.
Explain why a simpler alternative does not address the gap, and distinguish
requirements from implementation choices.

\begin{algorithm}[htbp]
  \caption{Illustrative candidate-selection protocol}
  \label{alg:selection}
  \begin{algorithmic}[1]
    \Require Candidate set $\Theta$, validation set $\mathcal{V}$, constraint $g$
    \Ensure Best feasible candidate, or explicit failure
    \State $\theta^{\star}\gets\bot$; $s^{\star}\gets+\infty$
    \For{each $\theta\in\Theta$}
      \If{$g(\theta,\mathcal{V})\leq 0$}
        \State $s\gets\operatorname{Evaluate}(\theta,\mathcal{V})$
        \If{$s<s^{\star}$}
          \State $(\theta^{\star},s^{\star})\gets(\theta,s)$
        \EndIf
      \EndIf
    \EndFor
    \State \Return $\theta^{\star}$
  \end{algorithmic}
\end{algorithm}

\autoref{alg:selection} permits explicit failure when no candidate passes.
State the budget, stopping condition, tie-breaking rule, and treatment of
invalid outputs. Selection uses validation data; reserve the held-out test
set for the final evaluation.

\subsection{Bound the claim}
\begin{proposition}[Finite-set selection]
If $\Theta$ is finite, all evaluations terminate, and the scores are finite
real numbers, \autoref{alg:selection} returns a minimum-score feasible
candidate whenever one exists.
\end{proposition}
\begin{proof}
After each iteration, the stored candidate has the smallest score among
the feasible candidates visited so far. The update preserves this invariant.
The finite loop visits every candidate, which establishes the result.
\end{proof}
This illustrative statement concerns only the enumerated candidate set.
It says nothing about unobserved inputs, global optimality, or deployment
security. State those limits alongside any theorem or empirical conclusion.
